\documentclass[aspectratio=169,10pt]{beamer}
\usetheme{Madrid}
\usecolortheme{default}
\usepackage[utf8]{inputenc}
\usepackage[T1]{fontenc}
\usepackage{lmodern}
\usepackage{amsmath,amssymb}
\usepackage{enumitem}
\usepackage{xcolor}

\setbeamertemplate{navigation symbols}{}
\setbeamertemplate{footline}[frame number]
\setbeamertemplate{itemize items}[circle]
\setbeamerfont{frametitle}{size=\large,series=\bfseries}

\title{CIA Part 1 -- Questions 60--118}
\subtitle{Governance, Risk, Controls, Ethics, IT, Fraud \& QAIP}
\author{CIA Study Prep}
\date{\today}

\begin{document}
	
	\begin{frame}
		\titlepage
	\end{frame}
	
	\begin{frame}{Instructions}
		\begin{itemize}
			\item Each question slide is followed by a separate answer/explanation slide.
			\item Questions are written at or above real exam difficulty.
			\item Distractors are designed to be plausible; read every word carefully.
		\end{itemize}
	\end{frame}
	
	% ============================================================
	% Q60
	% ============================================================
	\begin{frame}[t]{Q60}
		\small
		\textbf{Question:} The CAE of a listed company is preparing the annual audit plan. The audit committee has requested that internal audit cover "all significant risks." During planning, the CAE learns that management's risk register omits several risks related to a new digital lending platform (cyber, model risk, third-party data privacy), and that the platform is expected to generate 35\% of revenue within 18 months. The CAE also learns that the platform's vendor has not provided a SOC 2 report and that no internal audit work has been done on the platform. Which action is MOST consistent with the Standards?
		\begin{enumerate}[label=\Alph*.,itemsep=1pt,topsep=2pt]
			\item Rely on management's risk register because it is the approved risk document; schedule the platform for a future cycle.
			\item Include the platform risks in the audit universe, perform a risk assessment, prioritize coverage, and coordinate with other assurance providers.
			\item Defer planning until the vendor provides a SOC 2 report and management updates the risk register.
			\item Assign the platform audit to the IT vendor's internal auditors to leverage their expertise.
		\end{enumerate}
	\end{frame}
	
	\begin{frame}[t]{Q60 -- Answer \& Explanation}
		\small
		\textbf{Correct Answer: B} \\[4pt]
		\textbf{Why:} Standards 2010/2010.A1/2010.A2 require the CAE to consider the organization's risk profile, including emerging risks, and to develop a risk-based plan. A material new platform with unmitigated risks must be assessed and covered. \\[4pt]
		\textbf{A wrong:} The risk register is an input, not the sole basis; emerging risks must be considered. \\[2pt]
		\textbf{C wrong:} Waiting for the vendor or management delays assurance on a material risk. \\[2pt]
		\textbf{D wrong:} The vendor's auditors are not independent of the vendor and cannot provide objective assurance to the organization.
	\end{frame}
	
	% ============================================================
	% Q61
	% ============================================================
	\begin{frame}[t]{Q61}
		\small
		\textbf{Question:} During an engagement, an auditor discovers that the CFO has been directing large "consulting" payments to a company owned by a close friend, and that the CFO approved the related contracts. The auditor's draft report includes this finding. The CEO asks the auditor to "hold the finding" until the CFO's upcoming retirement, to "avoid embarrassment." Which action is MOST appropriate?
		\begin{enumerate}[label=\Alph*.,itemsep=1pt,topsep=2pt]
			\item Hold the finding as requested because the CFO is retiring soon.
			\item Report the finding promptly and escalate the CEO's request to the audit committee or board.
			\item Report the finding to the external auditor only.
			\item Remove the CFO's name from the finding but keep the finding.
		\end{enumerate}
	\end{frame}
	
	\begin{frame}[t]{Q61 -- Answer \& Explanation}
		\small
		\textbf{Correct Answer: B} \\[4pt]
		\textbf{Why:} Integrity/Objectivity (Code) and Standard 2600 require prompt reporting of significant issues, including potential fraud/conflict of interest, to the board/audit committee. Suppression is prohibited. \\[4pt]
		\textbf{A wrong:} Retirement timing does not justify delaying a significant finding. \\[2pt]
		\textbf{C wrong:} Reporting only to external audit does not satisfy board reporting. \\[2pt]
		\textbf{D wrong:} Removing the name does not cure the suppression or the substance.
	\end{frame}
	
	% ============================================================
	% Q62
	% ============================================================
	\begin{frame}[t]{Q62}
		\small
		\textbf{Question:} An auditor is testing the operating effectiveness of a control requiring two approvals for new vendor additions. The auditor selects a sample of 50 vendor additions and finds that 3 were approved by only one person. The auditor also finds that the system does not log approvals, so the audit trail is incomplete. Which conclusion is MOST appropriate?
		\begin{enumerate}[label=\Alph*.,itemsep=1pt,topsep=2pt]
			\item The control is effective because 47 of 50 were properly approved.
			\item The control is not operating effectively; deviations exist and the lack of an audit trail impairs evidence of operation.
			\item The control is effective because the deviations were "only 3."
			\item The auditor should accept the control because management approved the process.
		\end{enumerate}
	\end{frame}
	
	\begin{frame}[t]{Q62 -- Answer \& Explanation}
		\small
		\textbf{Correct Answer: B} \\[4pt]
		\textbf{Why:} Deviations indicate the control is not consistently operating, and the missing audit trail undermines evidence of operation (Standard 2300). This should be reported. \\[4pt]
		\textbf{A wrong:} Any deviation matters, especially for a key control. \\[2pt]
		\textbf{C wrong:} "Only 3" is not an acceptable rationale. \\[2pt]
		\textbf{D wrong:} Management approval of the process does not validate operation.
	\end{frame}
	
	% ============================================================
	% Q63
	% ============================================================
	\begin{frame}[t]{Q63}
		\small
		\textbf{Question:} The CAE is evaluating whether the internal audit activity's QAIP is effective. The activity has ongoing internal assessments, a periodic internal assessment performed last year, but the last external assessment was seven years ago and the CAE has not disclosed the nonconformance to the board. Which statement is MOST accurate?
		\begin{enumerate}[label=\Alph*.,itemsep=1pt,topsep=2pt]
			\item The QAIP is effective because internal assessments are current.
			\item The QAIP is nonconforming; Standard 1312 requires an external assessment at least every five years, and Standard 1320 requires disclosure of nonconformance to the board.
			\item The external assessment is optional if internal assessments are clean.
			\item The QAIP is effective because no significant issues were found internally.
		\end{enumerate}
	\end{frame}
	
	\begin{frame}[t]{Q63 -- Answer \& Explanation}
		\small
		\textbf{Correct Answer: B} \\[4pt]
		\textbf{Why:} Standard 1312 requires external assessment at least every five years; Standard 1320 requires reporting nonconformance to senior management and the board. \\[4pt]
		\textbf{A wrong:} Internal assessments do not replace external assessment. \\[2pt]
		\textbf{C wrong:} External assessment is mandatory. \\[2pt]
		\textbf{D wrong:} Clean internal assessments do not waive the requirement.
	\end{frame}
	
	% ============================================================
	% Q64
	% ============================================================
	\begin{frame}[t]{Q64}
		\small
		\textbf{Question:} An auditor is testing the completeness of recorded revenue. The auditor traces a sample of shipping documents to sales invoices and finds that 10\% of shipments were invoiced in the following period. Management states that revenue was recorded when the customer accepted the goods, not when shipped. The auditor reviews the sales contracts and finds that control transfers at shipment. Which conclusion is MOST appropriate?
		\begin{enumerate}[label=\Alph*.,itemsep=1pt,topsep=2pt]
			\item Revenue is complete because management's policy is consistent.
			\item Revenue is likely understated; control transferred at shipment, so revenue should have been recorded earlier.
			\item Revenue is overstated because invoices were issued late.
			\item The auditor should accept management's policy because it is documented.
		\end{enumerate}
	\end{frame}
	
	\begin{frame}[t]{Q64 -- Answer \& Explanation}
		\small
		\textbf{Correct Answer: B} \\[4pt]
		\textbf{Why:} Revenue recognition follows transfer of control. If contracts specify shipment, recording at customer acceptance understates revenue in the period. \\[4pt]
		\textbf{A wrong:} Consistency with a wrong policy does not make it correct. \\[2pt]
		\textbf{C wrong:} Late invoicing indicates understatement, not overstatement. \\[2pt]
		\textbf{D wrong:} Documentation does not override the contractual terms.
	\end{frame}
	
	% ============================================================
	% Q65
	% ============================================================
	\begin{frame}[t]{Q65}
		\small
		\textbf{Question:} The CAE is asked to assess the organization's governance of a joint venture in a foreign country. The joint venture has a local board, but the organization has only one representative, who is also the CFO, and the CFO has signing authority for all transactions. Which conclusion is MOST appropriate?
		\begin{enumerate}[label=\Alph*.,itemsep=1pt,topsep=2pt]
			\item Governance is adequate because the organization has a board representative.
			\item Governance is deficient; concentrating board representation and signing authority in one executive impairs oversight and creates a significant control risk.
			\item Joint venture governance is outside internal audit's scope.
			\item The auditor should ignore the issue because it is a foreign jurisdiction.
		\end{enumerate}
	\end{frame}
	
	\begin{frame}[t]{Q65 -- Answer \& Explanation}
		\small
		\textbf{Correct Answer: B} \\[4pt]
		\textbf{Why:} Effective governance requires independent oversight and segregation of duties. Concentrating board representation and signing authority in one executive (especially the CFO) is a significant governance/control weakness. \\[4pt]
		\textbf{A wrong:} A single representative is not adequate oversight. \\[2pt]
		\textbf{C wrong:} Joint ventures are within internal audit's scope. \\[2pt]
		\textbf{D wrong:} Foreign jurisdiction does not exempt governance review.
	\end{frame}
	
	% ============================================================
	% Q66
	% ============================================================
	\begin{frame}[t]{Q66}
		\small
		\textbf{Question:} An auditor is evaluating the organization's fraud risk assessment for the sales function. Management identifies "fictitious sales" as a risk but has no controls to address it and no monitoring. Which conclusion is MOST appropriate?
		\begin{enumerate}[label=\Alph*.,itemsep=1pt,topsep=2pt]
			\item The fraud risk assessment is adequate because the risk was identified.
			\item The fraud risk assessment is incomplete; identified risks require controls and monitoring to be effectively managed.
			\item Fraud risk assessment is not required for sales.
			\item The auditor should ignore the lack of controls because management identified the risk.
		\end{enumerate}
	\end{frame}
	
	\begin{frame}[t]{Q66 -- Answer \& Explanation}
		\small
		\textbf{Correct Answer: B} \\[4pt]
		\textbf{Why:} Identifying a risk without controls or monitoring does not manage it. The assessment should be linked to responses and monitoring. \\[4pt]
		\textbf{A wrong:} Identification alone is insufficient. \\[2pt]
		\textbf{C wrong:} Sales is a high-fraud-risk area requiring assessment. \\[2pt]
		\textbf{D wrong:} Lack of controls is a reportable weakness.
	\end{frame}
	
	% ============================================================
	% Q67
	% ============================================================
	\begin{frame}[t]{Q67}
		\small
		\textbf{Question:} The CAE is evaluating whether the organization's internal audit activity has sufficient authority. The charter grants unrestricted access to records, personnel, and property, and the audit committee approves the charter and plan. However, management has told staff not to cooperate with internal audit without management approval. Which conclusion is MOST appropriate?
		\begin{enumerate}[label=\Alph*.,itemsep=1pt,topsep=2pt]
			\item Authority is adequate because the charter grants unrestricted access.
			\item Authority is impaired in practice; management's instruction undermines the charter and internal audit's ability to perform engagements.
			\item Authority is adequate because the audit committee approves the plan.
			\item Authority is not relevant to internal audit.
		\end{enumerate}
	\end{frame}
	
	\begin{frame}[t]{Q67 -- Answer \& Explanation}
		\small
		\textbf{Correct Answer: B} \\[4pt]
		\textbf{Why:} Authority must be both granted (charter) and respected in practice. Management's instruction not to cooperate impairs internal audit's ability to perform and should be escalated. \\[4pt]
		\textbf{A wrong:} Charter authority is undermined by practice. \\[2pt]
		\textbf{C wrong:} Plan approval does not ensure cooperation. \\[2pt]
		\textbf{D wrong:} Authority is essential.
	\end{frame}
	
	% ============================================================
	% Q68
	% ============================================================
	\begin{frame}[t]{Q68}
		\small
		\textbf{Question:} An auditor is testing the organization's controls over the financial close. The auditor finds that the same individual prepares, reviews, and posts consolidation entries, and that the individual has system access to override the consolidation. Which conclusion is MOST appropriate?
		\begin{enumerate}[label=\Alph*.,itemsep=1pt,topsep=2pt]
			\item Controls are effective because the individual is experienced.
			\item Controls are deficient; incompatible duties and override access create a significant risk of undetected error/fraud.
			\item Controls are effective because consolidation is routine.
			\item The auditor should not assess the financial close.
		\end{enumerate}
	\end{frame}
	
	\begin{frame}[t]{Q68 -- Answer \& Explanation}
		\small
		\textbf{Correct Answer: B} \\[4pt]
		\textbf{Why:} Combining prepare/review/post and override access is a severe segregation-of-duties and override risk. \\[4pt]
		\textbf{A wrong:} Experience is not a control. \\[2pt]
		\textbf{C wrong:} Routine processes can still be manipulated. \\[2pt]
		\textbf{D wrong:} Financial close is a core audit area.
	\end{frame}
	
	% ============================================================
	% Q69
	% ============================================================
	\begin{frame}[t]{Q69}
		\small
		\textbf{Question:} The CAE is evaluating the organization's risk appetite for cybersecurity. The board has stated "zero tolerance for data breaches," but management has not funded basic security controls and has accepted known vulnerabilities. Which conclusion is MOST appropriate?
		\begin{enumerate}[label=\Alph*.,itemsep=1pt,topsep=2pt]
			\item Risk appetite and behavior are aligned because the board stated zero tolerance.
			\item There is a disconnect between stated risk appetite and actual behavior; the organization is not resourcing controls consistent with its appetite.
			\item Cybersecurity risk is not a governance issue.
			\item The auditor should ignore the disconnect.
		\end{enumerate}
	\end{frame}
	
	\begin{frame}[t]{Q69 -- Answer \& Explanation}
		\small
		\textbf{Correct Answer: B} \\[4pt]
		\textbf{Why:} Risk appetite must be supported by resources and actions. A stated zero tolerance without funding/controls indicates misalignment and a weak risk culture. \\[4pt]
		\textbf{A wrong:} Statements without action are not alignment. \\[2pt]
		\textbf{C wrong:} Cybersecurity is a governance risk. \\[2pt]
		\textbf{D wrong:} Misalignment is a reportable governance weakness.
	\end{frame}
	
	% ============================================================
	% Q70
	% ============================================================
	\begin{frame}[t]{Q70}
		\small
		\textbf{Question:} An auditor is assessing the organization's third-party risk management for a critical cloud provider. The provider has a SOC 2 Type II report, but the report covers a period ending 14 months ago and the provider has since changed its subservice providers. Which conclusion is MOST appropriate?
		\begin{enumerate}[label=\Alph*.,itemsep=1pt,topsep=2pt]
			\item The SOC 2 report is sufficient because it is a Type II report.
			\item The SOC 2 report is stale and does not cover recent subservice provider changes; the organization should obtain updated assurance and assess the changes.
			\item SOC 2 reports are irrelevant to internal audit.
			\item The organization should terminate the provider immediately.
		\end{enumerate}
	\end{frame}
	
	\begin{frame}[t]{Q70 -- Answer \& Explanation}
		\small
		\textbf{Correct Answer: B} \\[4pt]
		\textbf{Why:} SOC 2 Type II covers a period; a 14-month-old report and subsequent subservice changes reduce its relevance. Updated assurance is needed. \\[4pt]
		\textbf{A wrong:} Type II is useful but must be current and cover relevant changes. \\[2pt]
		\textbf{C wrong:} SOC reports are relevant evidence. \\[2pt]
		\textbf{D wrong:} Termination is an overreaction; remediation/updated assurance is appropriate.
	\end{frame}
	
	% ============================================================
	% Q71
	% ============================================================
	\begin{frame}[t]{Q71}
		\small
		\textbf{Question:} The CAE is asked to evaluate whether the organization's board has effective oversight of the audit function. The board receives audit reports but has no charter, no independent audit committee, and the CEO approves the CAE's performance evaluation. Which conclusion is MOST appropriate?
		\begin{enumerate}[label=\Alph*.,itemsep=1pt,topsep=2pt]
			\item Oversight is effective because the board receives audit reports.
			\item Oversight is deficient; lack of a charter, independent audit committee, and CEO-controlled CAE evaluation impair internal audit independence and board oversight.
			\item Oversight is effective because the CEO is accountable.
			\item Oversight is outside internal audit's scope.
		\end{enumerate}
	\end{frame}
	
	\begin{frame}[t]{Q71 -- Answer \& Explanation}
		\small
		\textbf{Correct Answer: B} \\[4pt]
		\textbf{Why:} Internal audit independence requires a board-approved charter, functional reporting to an independent audit committee, and CAE evaluation by the board. CEO-controlled evaluation impairs independence. \\[4pt]
		\textbf{A wrong:} Receiving reports is not effective oversight. \\[2pt]
		\textbf{C wrong:} CEO accountability does not ensure independence. \\[2pt]
		\textbf{D wrong:} Internal audit assesses governance.
	\end{frame}
	
	% ============================================================
	% Q72
	% ============================================================
	\begin{frame}[t]{Q72}
		\small
		\textbf{Question:} An auditor is testing the organization's controls over the use of artificial intelligence (AI) in credit decisions. The auditor finds that the AI model was developed by an external vendor, has not been validated, and that management cannot explain the model's logic. Which conclusion is MOST appropriate?
		\begin{enumerate}[label=\Alph*.,itemsep=1pt,topsep=2pt]
			\item Controls are adequate because the model was developed by a vendor.
			\item Controls are deficient; unvalidated AI models with unexplainable logic create significant model, compliance, and reputational risk.
			\item AI models do not require validation.
			\item The auditor should not assess AI models.
		\end{enumerate}
	\end{frame}
	
	\begin{frame}[t]{Q72 -- Answer \& Explanation}
		\small
		\textbf{Correct Answer: B} \\[4pt]
		\textbf{Why:} AI models used in credit decisions require validation, documentation, and explainability to manage model risk and regulatory compliance. \\[4pt]
		\textbf{A wrong:} Vendor development does not eliminate the need for validation. \\[2pt]
		\textbf{C wrong:} AI models require validation. \\[2pt]
		\textbf{D wrong:} AI risk is a core internal audit area.
	\end{frame}
	
	% ============================================================
	% Q73
	% ============================================================
	\begin{frame}[t]{Q73}
		\small
		\textbf{Question:} The CAE is evaluating the organization's fraud hotline. The hotline is managed by the internal audit function, reports go directly to the CAE, and the CAE investigates all reports. Which statement is MOST accurate?
		\begin{enumerate}[label=\Alph*.,itemsep=1pt,topsep=2pt]
			\item The hotline is effective because internal audit is independent.
			\item The hotline may create an independence/self-review issue if internal audit investigates reports about its own function; oversight by the audit committee is important.
			\item The hotline is ineffective because internal audit manages it.
			\item Hotlines should never be managed by internal audit.
		\end{enumerate}
	\end{frame}
	
	\begin{frame}[t]{Q73 -- Answer \& Explanation}
		\small
		\textbf{Correct Answer: B} \\[4pt]
		\textbf{Why:} Internal audit managing a hotline is common, but reports involving internal audit itself create a self-review threat. Audit committee oversight and escalation protocols are important safeguards. \\[4pt]
		\textbf{A wrong:} Independence alone does not resolve self-review for reports about internal audit. \\[2pt]
		\textbf{C wrong:} Internal audit management is not inherently ineffective. \\[2pt]
		\textbf{D wrong:} Internal audit can manage hotlines with safeguards.
	\end{frame}
	
	% ============================================================
	% Q74
	% ============================================================
	\begin{frame}[t]{Q74}
		\small
		\textbf{Question:} An auditor is testing the organization's controls over the financial close. The auditor finds that the controller prepares, reviews, and posts consolidation entries, and that the controller has system access to override the consolidation. Which conclusion is MOST appropriate?
		\begin{enumerate}[label=\Alph*.,itemsep=1pt,topsep=2pt]
			\item Controls are effective because the controller is experienced.
			\item Controls are deficient; incompatible duties and override access create a significant risk of undetected error/fraud.
			\item Controls are effective because consolidation is routine.
			\item The auditor should not assess the financial close.
		\end{enumerate}
	\end{frame}
	
	\begin{frame}[t]{Q74 -- Answer \& Explanation}
		\small
		\textbf{Correct Answer: B} \\[4pt]
		\textbf{Why:} Combining prepare/review/post and override access is a severe segregation-of-duties and override risk. \\[4pt]
		\textbf{A wrong:} Experience is not a control. \\[2pt]
		\textbf{B wrong:} Routine processes can still be manipulated. \\[2pt]
		\textbf{D wrong:} Financial close is a core audit area.
	\end{frame}
	
	% ============================================================
	% Q75
	% ============================================================
	\begin{frame}[t]{Q75}
		\small
		\textbf{Question:} The CAE is evaluating the organization's risk appetite for cybersecurity. The board has stated "zero tolerance for data breaches," but management has not funded basic security controls and has accepted known vulnerabilities. Which conclusion is MOST appropriate?
		\begin{enumerate}[label=\Alph*.,itemsep=1pt,topsep=2pt]
			\item Risk appetite and behavior are aligned because the board stated zero tolerance.
			\item There is a disconnect between stated risk appetite and actual behavior; the organization is not resourcing controls consistent with its appetite.
			\item Cybersecurity risk is not a governance issue.
			\item The auditor should ignore the disconnect.
		\end{enumerate}
	\end{frame}
	
	\begin{frame}[t]{Q75 -- Answer \& Explanation}
		\small
		\textbf{Correct Answer: B} \\[4pt]
		\textbf{Why:} Risk appetite must be supported by resources and actions. A stated zero tolerance without funding/controls indicates misalignment and a weak risk culture. \\[4pt]
		\textbf{A wrong:} Statements without action are not alignment. \\[2pt]
		\textbf{B wrong:} Cybersecurity is a governance risk. \\[2pt]
		\textbf{D wrong:} Misalignment is a reportable governance weakness.
	\end{frame}
	
	% ============================================================
	% Q76
	% ============================================================
	\begin{frame}[t]{Q76}
		\small
		\textbf{Question:} An auditor is testing the organization's controls over related-party transactions. The auditor finds that the CEO's spouse owns a consulting firm that received \$3 million in contracts, and that the transactions were not disclosed to the board or in the financial statements. Which conclusion is MOST appropriate?
		\begin{enumerate}[label=\Alph*.,itemsep=1pt,topsep=2pt]
			\item The transactions are acceptable because the firm is competitive.
			\item The transactions represent a significant control and disclosure deficiency; related-party transactions must be identified, approved, and disclosed.
			\item Related-party transactions are outside the scope of internal audit.
			\item The auditor should ignore the transactions because they involve a spouse.
		\end{enumerate}
	\end{frame}
	
	\begin{frame}[t]{Q76 -- Answer \& Explanation}
		\small
		\textbf{Correct Answer: B} \\[4pt]
		\textbf{Why:} Related-party transactions require identification, independent approval, and disclosure. Undisclosed transactions with the CEO's spouse's firm are a significant governance/control failure. \\[4pt]
		\textbf{A wrong:} Competitiveness does not eliminate disclosure requirements. \\[2pt]
		\textbf{C wrong:} Related-party transactions are a core audit area. \\[2pt]
		\textbf{D wrong:} Spousal relationships can be fraud indicators.
	\end{frame}
	
	% ============================================================
	% Q77
	% ============================================================
	\begin{frame}[t]{Q77}
		\small
		\textbf{Question:} The CAE is evaluating the organization's cybersecurity governance. The organization has no board-level cybersecurity expertise, no regular cybersecurity reporting to the board, and no incident response plan. Which conclusion is MOST appropriate?
		\begin{enumerate}[label=\Alph*.,itemsep=1pt,topsep=2pt]
			\item Cybersecurity governance is adequate because IT handles security.
			\item Cybersecurity governance is deficient; board oversight, regular reporting, and an incident response plan are essential.
			\item Cybersecurity is purely an IT issue, not a governance issue.
			\item The auditor should not assess cybersecurity.
		\end{enumerate}
	\end{frame}
	
	\begin{frame}[t]{Q77 -- Answer \& Explanation}
		\small
		\textbf{Correct Answer: B} \\[4pt]
		\textbf{Why:} Cybersecurity is a board-level governance risk. Lack of expertise, reporting, and an incident response plan are significant governance deficiencies. \\[4pt]
		\textbf{A wrong:} Security is not solely an IT issue; governance oversight is required. \\[2pt]
		\textbf{C wrong:} Cybersecurity is a governance and enterprise risk. \\[2pt]
		\textbf{D wrong:} Internal audit assesses cybersecurity governance.
	\end{frame}
	
	% ============================================================
	% Q78
	% ============================================================
	\begin{frame}[t]{Q78}
		\small
		\textbf{Question:} An auditor is evaluating the organization's fraud risk assessment. Management's assessment is based solely on a brainstorming session with the CFO and controller, with no input from operations, sales, or IT, and no consideration of fraud schemes specific to the industry. Which conclusion is MOST appropriate?
		\begin{enumerate}[label=\Alph*.,itemsep=1pt,topsep=2pt]
			\item The fraud risk assessment is adequate because senior finance staff participated.
			\item The fraud risk assessment is deficient; it lacks cross-functional input and industry-specific fraud schemes.
			\item The fraud risk assessment is adequate because brainstorming is a recognized technique.
			\item Fraud risk assessment is not an internal audit concern.
		\end{enumerate}
	\end{frame}
	
	\begin{frame}[t]{Q78 -- Answer \& Explanation}
		\small
		\textbf{Correct Answer: B} \\[4pt]
		\textbf{Why:} Effective fraud risk assessment requires cross-functional input and consideration of industry-specific schemes. A finance-only brainstorm is insufficient. \\[4pt]
		\textbf{A wrong:} Finance-only input is too narrow. \\[2pt]
		\textbf{C wrong:} Brainstorming is useful but must be broad and informed. \\[2pt]
		\textbf{D wrong:} Internal audit assesses fraud risk (Standard 1210.A2).
	\end{frame}
	
	% ============================================================
	% Q79
	% ============================================================
	\begin{frame}[t]{Q79}
		\small
		\textbf{Question:} The CAE is assessing the organization's internal control framework. The organization has a strong control environment and strong control activities but no monitoring activities and no information/communication process. Which conclusion is MOST appropriate?
		\begin{enumerate}[label=\Alph*.,itemsep=1pt,topsep=2pt]
			\item The internal control framework is effective because the control environment and activities are strong.
			\item The internal control framework is incomplete; all five COSO components (including monitoring and information/communication) are necessary.
			\item Monitoring and information/communication are optional.
			\item Only control activities matter.
		\end{enumerate}
	\end{frame}
	
	\begin{frame}[t]{Q79 -- Answer \& Explanation}
		\small
		\textbf{Correct Answer: B} \\[4pt]
		\textbf{Why:} COSO requires all five components: control environment, risk assessment, control activities, information/communication, and monitoring. Missing components make the framework incomplete. \\[4pt]
		\textbf{A wrong:} Two components do not make a complete framework. \\[2pt]
		\textbf{C wrong:} All components are required. \\[2pt]
		\textbf{D wrong:} Control activities alone are insufficient.
	\end{frame}
	
	% ============================================================
	% Q80
	% ============================================================
	\begin{frame}[t]{Q80}
		\small
		\textbf{Question:} An auditor is testing the organization's controls over the purchase of consulting services. The auditor finds that the same manager selects the consultant, negotiates the contract, approves the invoice, and confirms receipt of services, with no independent review. Which conclusion is MOST appropriate?
		\begin{enumerate}[label=\Alph*.,itemsep=1pt,topsep=2pt]
			\item Controls are effective because the manager is accountable for results.
			\item Controls are deficient; the manager has incompatible duties (selection, negotiation, approval, receipt) with no independent review.
			\item Controls are effective because consulting services are intangible.
			\item The auditor should not assess consulting purchases.
		\end{enumerate}
	\end{frame}
	
	\begin{frame}[t]{Q80 -- Answer \& Explanation}
		\small
		\textbf{Correct Answer: B} \\[4pt]
		\textbf{Why:} Combining selection, negotiation, approval, and receipt of services creates a significant segregation-of-duties weakness and fraud risk. \\[4pt]
		\textbf{A wrong:} Accountability is not a control. \\[2pt]
		\textbf{C wrong:} Intangibility does not reduce the need for controls. \\[2pt]
		\textbf{D wrong:} Consulting purchases are a legitimate audit area.
	\end{frame}
	
	% ============================================================
	% Q81
	% ============================================================
	\begin{frame}[t]{Q81}
		\small
		\textbf{Question:} The CAE is evaluating the organization's risk response to a significant cybersecurity risk. Management has purchased cyber insurance but has not implemented basic security controls (patching, access management, monitoring). Which conclusion is MOST appropriate?
		\begin{enumerate}[label=\Alph*.,itemsep=1pt,topsep=2pt]
			\item The risk response is adequate because insurance transfers the risk.
			\item The risk response is inadequate; insurance transfers financial impact but does not reduce the likelihood of a breach, and basic controls are missing.
			\item Insurance eliminates the need for controls.
			\item Cybersecurity risk is not an internal audit concern.
		\end{enumerate}
	\end{frame}
	
	\begin{frame}[t]{Q81 -- Answer \& Explanation}
		\small
		\textbf{Correct Answer: B} \\[4pt]
		\textbf{Why:} Insurance is a risk-sharing response, not a substitute for preventive/detective controls. Basic security controls are essential to reduce likelihood. \\[4pt]
		\textbf{A wrong:} Insurance does not prevent breaches. \\[2pt]
		\textbf{C wrong:} Insurance does not eliminate the need for controls. \\[2pt]
		\textbf{D wrong:} Cybersecurity is a core internal audit area.
	\end{frame}
	
	% ============================================================
	% Q82
	% ============================================================
	\begin{frame}[t]{Q82}
		\small
		\textbf{Question:} An auditor is assessing the organization's use of a whistleblower hotline. The hotline is outsourced to an independent provider, reports go directly to the audit committee, and the audit committee reviews all reports. However, no reports have been filed in two years, and employee surveys indicate fear of retaliation. Which conclusion is MOST appropriate?
		\begin{enumerate}[label=\Alph*.,itemsep=1pt,topsep=2pt]
			\item The hotline is effective because it is independent and reports to the audit committee.
			\item The hotline may be underutilized due to fear of retaliation; the organization should address the retaliation risk and promote the hotline.
			\item Zero reports prove the organization is ethical.
			\item The auditor should ignore employee surveys.
		\end{enumerate}
	\end{frame}
	
	\begin{frame}[t]{Q82 -- Answer \& Explanation}
		\small
		\textbf{Correct Answer: B} \\[4pt]
		\textbf{Why:} A hotline's effectiveness depends on trust and non-retaliation. Zero reports combined with fear of retaliation suggests underutilization and a culture problem. \\[4pt]
		\textbf{A wrong:} Structure alone is insufficient. \\[2pt]
		\textbf{C wrong:} Zero reports may indicate fear, not ethics. \\[2pt]
		\textbf{D wrong:} Surveys provide relevant culture evidence.
	\end{frame}
	
	% ============================================================
	% Q83
	% ============================================================
	\begin{frame}[t]{Q83}
		\small
		\textbf{Question:} The CAE is asked to evaluate whether the organization's board has effective oversight of risk. The board receives a quarterly risk report that is prepared by management and contains only risk status ("green") with no details on emerging risks or risk appetite breaches. Which conclusion is MOST appropriate?
		\begin{enumerate}[label=\Alph*.,itemsep=1pt,topsep=2pt]
			\item Board risk oversight is effective because the board receives a quarterly report.
			\item Board risk oversight is likely ineffective; the report lacks detail, emerging risks, and appetite breach information.
			\item Risk reporting is management's responsibility, not the board's.
			\item The auditor should not assess board risk oversight.
		\end{enumerate}
	\end{frame}
	
	\begin{frame}[t]{Q83 -- Answer \& Explanation}
		\small
		\textbf{Correct Answer: B} \\[4pt]
		\textbf{Why:} Effective board risk oversight requires transparent, detailed reporting on emerging risks and appetite breaches. A "green-only" summary impairs oversight. \\[4pt]
		\textbf{A wrong:} Receiving a report is not the same as effective oversight. \\[2pt]
		\textbf{C wrong:} The board is responsible for oversight. \\[2pt]
		\textbf{D wrong:} Internal audit assesses governance/risk oversight.
	\end{frame}
	
	% ============================================================
	% Q84
	% ============================================================
	\begin{frame}[t]{Q84}
		\small
		\textbf{Question:} An auditor is testing the organization's controls over the sales commission process. The auditor finds that sales managers can adjust commission rates and approve their own adjustments, with no independent review. Which conclusion is MOST appropriate?
		\begin{enumerate}[label=\Alph*.,itemsep=1pt,topsep=2pt]
			\item Controls are effective because sales managers understand commissions.
			\item Controls are deficient; sales managers can adjust rates and approve their own changes, creating a significant fraud risk.
			\item Controls are effective because commission adjustments are rare.
			\item The auditor should not assess sales commissions.
		\end{enumerate}
	\end{frame}
	
	\begin{frame}[t]{Q84 -- Answer \& Explanation}
		\small
		\textbf{Correct Answer: B} \\[4pt]
		\textbf{Why:} The ability to adjust rates and self-approve creates a significant fraud risk (self-dealing). Independent review is required. \\[4pt]
		\textbf{A wrong:} Knowledge is not a control. \\[2pt]
		\textbf{C wrong:} Rarity does not reduce the fraud risk. \\[2pt]
		\textbf{D wrong:} Commissions are a core audit area.
	\end{frame}
	
	% ============================================================
	% Q85
	% ============================================================
	\begin{frame}[t]{Q85}
		\small
		\textbf{Question:} The CAE is evaluating the organization's compliance with its own code of conduct. The CAE finds that senior executives routinely accept expensive gifts from vendors and that no one has been disciplined. Which conclusion is MOST appropriate?
		\begin{enumerate}[label=\Alph*.,itemsep=1pt,topsep=2pt]
			\item The code of conduct is effective because it is written.
			\item The code of conduct is not effective; senior executive behavior and lack of enforcement undermine the code.
			\item Gifts are acceptable if they are expensive.
			\item The auditor should ignore executive behavior.
		\end{enumerate}
	\end{frame}
	
	\begin{frame}[t]{Q85 -- Answer \& Explanation}
		\small
		\textbf{Correct Answer: B} \\[4pt]
		\textbf{Why:} Tone at the top and enforcement are essential. Senior executives accepting gifts without consequences undermine the code's effectiveness. \\[4pt]
		\textbf{A wrong:} A written code without enforcement is ineffective. \\[2pt]
		\textbf{C wrong:} Expensive gifts create significant conflicts. \\[2pt]
		\textbf{D wrong:} Executive behavior is a key governance indicator.
	\end{frame}
	
	% ============================================================
	% Q86
	% ============================================================
	\begin{frame}[t]{Q86}
		\small
		\textbf{Question:} An auditor is assessing the organization's controls over the use of spreadsheets in financial reporting. The auditor finds that a critical revenue calculation is performed in a spreadsheet with no version control, no formula documentation, and no independent review. Which conclusion is MOST appropriate?
		\begin{enumerate}[label=\Alph*.,itemsep=1pt,topsep=2pt]
			\item Controls are adequate because spreadsheets are common tools.
			\item Controls are deficient; critical spreadsheets require version control, documentation, and independent review.
			\item Spreadsheets are always inappropriate for financial reporting.
			\item The auditor should not assess spreadsheet controls.
		\end{enumerate}
	\end{frame}
	
	\begin{frame}[t]{Q86 -- Answer \& Explanation}
		\small
		\textbf{Correct Answer: B} \\[4pt]
		\textbf{Why:} Critical spreadsheets are high-risk end-user computing tools requiring version control, documentation, and review to prevent errors/fraud. \\[4pt]
		\textbf{A wrong:} Common use does not reduce risk. \\[2pt]
		\textbf{C wrong:} Spreadsheets can be used with adequate controls. \\[2pt]
		\textbf{D wrong:} Spreadsheet controls are a valid audit area.
	\end{frame}
	
	% ============================================================
	% Q87
	% ============================================================
	\begin{frame}[t]{Q87}
		\small
		\textbf{Question:} The CAE is evaluating the organization's risk management for a major IT system implementation. The project has no risk register, no steering committee, and no independent oversight; the project manager reports only to the CIO. Which conclusion is MOST appropriate?
		\begin{enumerate}[label=\Alph*.,itemsep=1pt,topsep=2pt]
			\item Project risk management is adequate because the CIO is accountable.
			\item Project risk management is deficient; major IT projects require a risk register, steering committee, and independent oversight.
			\item IT projects do not require risk management.
			\item The auditor should not assess IT project risk.
		\end{enumerate}
	\end{frame}
	
	\begin{frame}[t]{Q87 -- Answer \& Explanation}
		\small
		\textbf{Correct Answer: B} \\[4pt]
		\textbf{Why:} Major IT projects require structured risk management, governance (steering committee), and independent oversight to manage cost, schedule, and risk. \\[4pt]
		\textbf{A wrong:} CIO accountability alone is insufficient. \\[2pt]
		\textbf{C wrong:} All projects require risk management. \\[2pt]
		\textbf{D wrong:} IT project risk is a core audit area.
	\end{frame}
	
	% ============================================================
	% Q88
	% ============================================================
	\begin{frame}[t]{Q88}
		\small
		\textbf{Question:} An auditor is testing the organization's controls over journal entries. The auditor finds that 20 large manual entries were posted at year-end by the controller with descriptions such as "reclass" and "true-up," and that supporting documentation cannot be located. Which action is MOST appropriate?
		\begin{enumerate}[label=\Alph*.,itemsep=1pt,topsep=2pt]
			\item Accept the entries because the controller is authorized.
			\item Investigate the entries, request support, and report the lack of documentation as a significant deficiency.
			\item Delete the entries from the system.
			\item Ignore the entries because they are year-end adjustments.
		\end{enumerate}
	\end{frame}
	
	\begin{frame}[t]{Q88 -- Answer \& Explanation}
		\small
		\textbf{Correct Answer: B} \\[4pt]
		\textbf{Why:} Large manual year-end entries by senior management are a high fraud risk area. Lack of support is a significant deficiency requiring investigation and reporting. \\[4pt]
		\textbf{A wrong:} Authorization does not eliminate the need for support. \\[2pt]
		\textbf{C wrong:} Internal audit should not alter records. \\[2pt]
		\textbf{D wrong:} Year-end adjustments require documentation.
	\end{frame}
	
	% ============================================================
	% Q89
	% ============================================================
	\begin{frame}[t]{Q89}
		\small
		\textbf{Question:} The CAE is evaluating the organization's use of key risk indicators (KRIs). The organization tracks revenue, profit, and customer satisfaction but has no KRIs for compliance, cybersecurity, or supply chain risk. Which conclusion is MOST appropriate?
		\begin{enumerate}[label=\Alph*.,itemsep=1pt,topsep=2pt]
			\item KRI coverage is adequate because financial and customer metrics are tracked.
			\item KRI coverage is incomplete; the organization lacks KRIs for significant risk areas (compliance, cybersecurity, supply chain).
			\item KRIs are optional for all organizations.
			\item The auditor should not assess KRIs.
		\end{enumerate}
	\end{frame}
	
	\begin{frame}[t]{Q89 -- Answer \& Explanation}
		\small
		\textbf{Correct Answer: B} \\[4pt]
		\textbf{Why:} KRIs should cover the organization's significant risks. Omitting compliance, cybersecurity, and supply chain KRIs leaves key risks unmonitored. \\[4pt]
		\textbf{A wrong:} Financial/customer metrics alone are insufficient. \\[2pt]
		\textbf{C wrong:} KRIs are an important monitoring tool. \\[2pt]
		\textbf{D wrong:} Internal audit assesses risk monitoring.
	\end{frame}
	
	% ============================================================
	% Q90
	% ============================================================
	\begin{frame}[t]{Q90}
		\small
		\textbf{Question:} An auditor is evaluating the organization's controls over the financial close. The auditor finds that the same individual prepares, reviews, and posts consolidation entries, and that the individual has system access to override the consolidation. Which conclusion is MOST appropriate?
		\begin{enumerate}[label=\Alph*.,itemsep=1pt,topsep=2pt]
			\item Controls are effective because the individual is experienced.
			\item Controls are deficient; incompatible duties and override access create a significant risk of undetected error/fraud.
			\item Controls are effective because consolidation is routine.
			\item The auditor should not assess the financial close.
		\end{enumerate}
	\end{frame}
	
	\begin{frame}[t]{Q90 -- Answer \& Explanation}
		\small
		\textbf{Correct Answer: B} \\[4pt]
		\textbf{Why:} Combining prepare/review/post and override access is a severe segregation-of-duties and override risk. \\[4pt]
		\textbf{A wrong:} Experience is not a control. \\[2pt]
		\textbf{C wrong:} Routine processes can still be manipulated. \\[2pt]
		\textbf{D wrong:} Financial close is a core audit area.
	\end{frame}
	
	% ============================================================
	% Q91
	% ============================================================
	\begin{frame}[t]{Q91}
		\small
		\textbf{Question:} The CAE is evaluating the organization's risk appetite for cybersecurity. The board has stated "zero tolerance for data breaches," but management has not funded basic security controls and has accepted known vulnerabilities. Which conclusion is MOST appropriate?
		\begin{enumerate}[label=\Alph*.,itemsep=1pt,topsep=2pt]
			\item Risk appetite and behavior are aligned because the board stated zero tolerance.
			\item There is a disconnect between stated risk appetite and actual behavior; the organization is not resourcing controls consistent with its appetite.
			\item Cybersecurity risk is not a governance issue.
			\item The auditor should ignore the disconnect.
		\end{enumerate}
	\end{frame}
	
	\begin{frame}[t]{Q91 -- Answer \& Explanation}
		\small
		\textbf{Correct Answer: B} \\[4pt]
		\textbf{Why:} Risk appetite must be supported by resources and actions. A stated zero tolerance without funding/controls indicates misalignment and a weak risk culture. \\[4pt]
		\textbf{A wrong:} Statements without action are not alignment. \\[2pt]
		\textbf{C wrong:} Cybersecurity is a governance risk. \\[2pt]
		\textbf{D wrong:} Misalignment is a reportable governance weakness.
	\end{frame}
	
	% ============================================================
	% Q92
	% ============================================================
	\begin{frame}[t]{Q92}
		\small
		\textbf{Question:} An auditor is testing the organization's controls over the use of artificial intelligence (AI) in credit decisions. The auditor finds that the AI model was developed by an external vendor, has not been validated, and that management cannot explain the model's logic. Which conclusion is MOST appropriate?
		\begin{enumerate}[label=\Alph*.,itemsep=1pt,topsep=2pt]
			\item Controls are adequate because the model was developed by a vendor.
			\item Controls are deficient; unvalidated AI models with unexplainable logic create significant model, compliance, and reputational risk.
			\item AI models do not require validation.
			\item The auditor should not assess AI models.
		\end{enumerate}
	\end{frame}
	
	\begin{frame}[t]{Q92 -- Answer \& Explanation}
		\small
		\textbf{Correct Answer: B} \\[4pt]
		\textbf{Why:} AI models used in credit decisions require validation, documentation, and explainability to manage model risk and regulatory compliance. \\[4pt]
		\textbf{A wrong:} Vendor development does not eliminate the need for validation. \\[2pt]
		\textbf{C wrong:} AI models require validation. \\[2pt]
		\textbf{D wrong:} AI risk is a core internal audit area.
	\end{frame}
	
	% ============================================================
	% Q93
	% ============================================================
	\begin{frame}[t]{Q93}
		\small
		\textbf{Question:} The CAE is evaluating the organization's fraud hotline. The hotline is managed by the internal audit function, reports go directly to the CAE, and the CAE investigates all reports. Which statement is MOST accurate?
		\begin{enumerate}[label=\Alph*.,itemsep=1pt,topsep=2pt]
			\item The hotline is effective because internal audit is independent.
			\item The hotline may create an independence/self-review issue if internal audit investigates reports about its own function; oversight by the audit committee is important.
			\item The hotline is ineffective because internal audit manages it.
			\item Hotlines should never be managed by internal audit.
		\end{enumerate}
	\end{frame}
	
	\begin{frame}[t]{Q93 -- Answer \& Explanation}
		\small
		\textbf{Correct Answer: B} \\[4pt]
		\textbf{Why:} Internal audit managing a hotline is common, but reports involving internal audit itself create a self-review threat. Audit committee oversight and escalation protocols are important safeguards. \\[4pt]
		\textbf{A wrong:} Independence alone does not resolve self-review for reports about internal audit. \\[2pt]
		\textbf{C wrong:} Internal audit management is not inherently ineffective. \\[2pt]
		\textbf{D wrong:} Internal audit can manage hotlines with safeguards.
	\end{frame}
	
	% ============================================================
	% Q94
	% ============================================================
	\begin{frame}[t]{Q94}
		\small
		\textbf{Question:} An auditor is testing the organization's controls over the financial close. The auditor finds that the controller prepares, reviews, and posts consolidation entries, and that the controller has system access to override the consolidation. Which conclusion is MOST appropriate?
		\begin{enumerate}[label=\Alph*.,itemsep=1pt,topsep=2pt]
			\item Controls are effective because the controller is experienced.
			\item Controls are deficient; incompatible duties and override access create a significant risk of undetected error/fraud.
			\item Controls are effective because consolidation is routine.
			\item The auditor should not assess the financial close.
		\end{enumerate}
	\end{frame}
	
	\begin{frame}[t]{Q94 -- Answer \& Explanation}
		\small
		\textbf{Correct Answer: B} \\[4pt]
		\textbf{Why:} Combining prepare/review/post and override access is a severe segregation-of-duties and override risk. \\[4pt]
		\textbf{A wrong:} Experience is not a control. \\[2pt]
		\textbf{C wrong:} Routine processes can still be manipulated. \\[2pt]
		\textbf{D wrong:} Financial close is a core audit area.
	\end{frame}
	
	% ============================================================
	% Q95
	% ============================================================
	\begin{frame}[t]{Q95}
		\small
		\textbf{Question:} The CAE is evaluating the organization's risk appetite for cybersecurity. The board has stated "zero tolerance for data breaches," but management has not funded basic security controls and has accepted known vulnerabilities. Which conclusion is MOST appropriate?
		\begin{enumerate}[label=\Alph*.,itemsep=1pt,topsep=2pt]
			\item Risk appetite and behavior are aligned because the board stated zero tolerance.
			\item There is a disconnect between stated risk appetite and actual behavior; the organization is not resourcing controls consistent with its appetite.
			\item Cybersecurity risk is not a governance issue.
			\item The auditor should ignore the disconnect.
		\end{enumerate}
	\end{frame}
	
	\begin{frame}[t]{Q95 -- Answer \& Explanation}
		\small
		\textbf{Correct Answer: B} \\[4pt]
		\textbf{Why:} Risk appetite must be supported by resources and actions. A stated zero tolerance without funding/controls indicates misalignment and a weak risk culture. \\[4pt]
		\textbf{A wrong:} Statements without action are not alignment. \\[2pt]
		\textbf{C wrong:} Cybersecurity is a governance risk. \\[2pt]
		\textbf{D wrong:} Misalignment is a reportable governance weakness.
	\end{frame}
	
	% ============================================================
	% Q96
	% ============================================================
	\begin{frame}[t]{Q96}
		\small
		\textbf{Question:} An auditor is testing the organization's controls over related-party transactions. The auditor finds that the CEO's spouse owns a consulting firm that received \$3 million in contracts, and that the transactions were not disclosed to the board or in the financial statements. Which conclusion is MOST appropriate?
		\begin{enumerate}[label=\Alph*.,itemsep=1pt,topsep=2pt]
			\item The transactions are acceptable because the firm is competitive.
			\item The transactions represent a significant control and disclosure deficiency; related-party transactions must be identified, approved, and disclosed.
			\item Related-party transactions are outside the scope of internal audit.
			\item The auditor should ignore the transactions because they involve a spouse.
		\end{enumerate}
	\end{frame}
	
	\begin{frame}[t]{Q96 -- Answer \& Explanation}
		\small
		\textbf{Correct Answer: B} \\[4pt]
		\textbf{Why:} Related-party transactions require identification, independent approval, and disclosure. Undisclosed transactions with the CEO's spouse's firm are a significant governance/control failure. \\[4pt]
		\textbf{A wrong:} Competitiveness does not eliminate disclosure requirements. \\[2pt]
		\textbf{C wrong:} Related-party transactions are a core audit area. \\[2pt]
		\textbf{D wrong:} Spousal relationships can be fraud indicators.
	\end{frame}
	
	% ============================================================
	% Q97
	% ============================================================
	\begin{frame}[t]{Q97}
		\small
		\textbf{Question:} The CAE is evaluating the organization's cybersecurity governance. The organization has no board-level cybersecurity expertise, no regular cybersecurity reporting to the board, and no incident response plan. Which conclusion is MOST appropriate?
		\begin{enumerate}[label=\Alph*.,itemsep=1pt,topsep=2pt]
			\item Cybersecurity governance is adequate because IT handles security.
			\item Cybersecurity governance is deficient; board oversight, regular reporting, and an incident response plan are essential.
			\item Cybersecurity is purely an IT issue, not a governance issue.
			\item The auditor should not assess cybersecurity.
		\end{enumerate}
	\end{frame}
	
	\begin{frame}[t]{Q97 -- Answer \& Explanation}
		\small
		\textbf{Correct Answer: B} \\[4pt]
		\textbf{Why:} Cybersecurity is a board-level governance risk. Lack of expertise, reporting, and an incident response plan are significant governance deficiencies. \\[4pt]
		\textbf{A wrong:} Security is not solely an IT issue; governance oversight is required. \\[2pt]
		\textbf{C wrong:} Cybersecurity is a governance and enterprise risk. \\[2pt]
		\textbf{D wrong:} Internal audit assesses cybersecurity governance.
	\end{frame}
	
	% ============================================================
	% Q98
	% ============================================================
	\begin{frame}[t]{Q98}
		\small
		\textbf{Question:} An auditor is evaluating the organization's fraud risk assessment. Management's assessment is based solely on a brainstorming session with the CFO and controller, with no input from operations, sales, or IT, and no consideration of fraud schemes specific to the industry. Which conclusion is MOST appropriate?
		\begin{enumerate}[label=\Alph*.,itemsep=1pt,topsep=2pt]
			\item The fraud risk assessment is adequate because senior finance staff participated.
			\item The fraud risk assessment is deficient; it lacks cross-functional input and industry-specific fraud schemes.
			\item The fraud risk assessment is adequate because brainstorming is a recognized technique.
			\item Fraud risk assessment is not an internal audit concern.
		\end{enumerate}
	\end{frame}
	
	\begin{frame}[t]{Q98 -- Answer \& Explanation}
		\small
		\textbf{Correct Answer: B} \\[4pt]
		\textbf{Why:} Effective fraud risk assessment requires cross-functional input and consideration of industry-specific schemes. A finance-only brainstorm is insufficient. \\[4pt]
		\textbf{A wrong:} Finance-only input is too narrow. \\[2pt]
		\textbf{C wrong:} Brainstorming is useful but must be broad and informed. \\[2pt]
		\textbf{D wrong:} Internal audit assesses fraud risk (Standard 1210.A2).
	\end{frame}
	
	% ============================================================
	% Q99
	% ============================================================
	\begin{frame}[t]{Q99}
		\small
		\textbf{Question:} The CAE is assessing the organization's internal control framework. The organization has a strong control environment and strong control activities but no monitoring activities and no information/communication process. Which conclusion is MOST appropriate?
		\begin{enumerate}[label=\Alph*.,itemsep=1pt,topsep=2pt]
			\item The internal control framework is effective because the control environment and activities are strong.
			\item The internal control framework is incomplete; all five COSO components (including monitoring and information/communication) are necessary.
			\item Monitoring and information/communication are optional.
			\item Only control activities matter.
		\end{enumerate}
	\end{frame}
	
	\begin{frame}[t]{Q99 -- Answer \& Explanation}
		\small
		\textbf{Correct Answer: B} \\[4pt]
		\textbf{Why:} COSO requires all five components: control environment, risk assessment, control activities, information/communication, and monitoring. Missing components make the framework incomplete. \\[4pt]
		\textbf{A wrong:} Two components do not make a complete framework. \\[2pt]
		\textbf{C wrong:} All components are required. \\[2pt]
		\textbf{D wrong:} Control activities alone are insufficient.
	\end{frame}
	
	% ============================================================
	% Q100
	% ============================================================
	\begin{frame}[t]{Q100}
		\small
		\textbf{Question:} An auditor is testing the organization's controls over the purchase of consulting services. The auditor finds that the same manager selects the consultant, negotiates the contract, approves the invoice, and confirms receipt of services, with no independent review. Which conclusion is MOST appropriate?
		\begin{enumerate}[label=\Alph*.,itemsep=1pt,topsep=2pt]
			\item Controls are effective because the manager is accountable for results.
			\item Controls are deficient; the manager has incompatible duties (selection, negotiation, approval, receipt) with no independent review.
			\item Controls are effective because consulting services are intangible.
			\item The auditor should not assess consulting purchases.
		\end{enumerate}
	\end{frame}
	
	\begin{frame}[t]{Q100 -- Answer \& Explanation}
		\small
		\textbf{Correct Answer: B} \\[4pt]
		\textbf{Why:} Combining selection, negotiation, approval, and receipt of services creates a significant segregation-of-duties weakness and fraud risk. \\[4pt]
		\textbf{A wrong:} Accountability is not a control. \\[2pt]
		\textbf{C wrong:} Intangibility does not reduce the need for controls. \\[2pt]
		\textbf{D wrong:} Consulting purchases are a legitimate audit area.
	\end{frame}
	
	% ============================================================
	% Q101
	% ============================================================
	\begin{frame}[t]{Q101}
		\small
		\textbf{Question:} The CAE is evaluating the organization's risk response to a significant cybersecurity risk. Management has purchased cyber insurance but has not implemented basic security controls (patching, access management, monitoring). Which conclusion is MOST appropriate?
		\begin{enumerate}[label=\Alph*.,itemsep=1pt,topsep=2pt]
			\item The risk response is adequate because insurance transfers the risk.
			\item The risk response is inadequate; insurance transfers financial impact but does not reduce the likelihood of a breach, and basic controls are missing.
			\item Insurance eliminates the need for controls.
			\item Cybersecurity risk is not an internal audit concern.
		\end{enumerate}
	\end{frame}
	
	\begin{frame}[t]{Q101 -- Answer \& Explanation}
		\small
		\textbf{Correct Answer: B} \\[4pt]
		\textbf{Why:} Insurance is a risk-sharing response, not a substitute for preventive/detective controls. Basic security controls are essential to reduce likelihood. \\[4pt]
		\textbf{A wrong:} Insurance does not prevent breaches. \\[2pt]
		\textbf{C wrong:} Insurance does not eliminate the need for controls. \\[2pt]
		\textbf{D wrong:} Cybersecurity is a core internal audit area.
	\end{frame}
	
	% ============================================================
	% Q102
	% ============================================================
	\begin{frame}[t]{Q102}
		\small
		\textbf{Question:} An auditor is assessing the organization's use of a whistleblower hotline. The hotline is outsourced to an independent provider, reports go directly to the audit committee, and the audit committee reviews all reports. However, no reports have been filed in two years, and employee surveys indicate fear of retaliation. Which conclusion is MOST appropriate?
		\begin{enumerate}[label=\Alph*.,itemsep=1pt,topsep=2pt]
			\item The hotline is effective because it is independent and reports to the audit committee.
			\item The hotline may be underutilized due to fear of retaliation; the organization should address the retaliation risk and promote the hotline.
			\item Zero reports prove the organization is ethical.
			\item The auditor should ignore employee surveys.
		\end{enumerate}
	\end{frame}
	
	\begin{frame}[t]{Q102 -- Answer \& Explanation}
		\small
		\textbf{Correct Answer: B} \\[4pt]
		\textbf{Why:} A hotline's effectiveness depends on trust and non-retaliation. Zero reports combined with fear of retaliation suggests underutilization and a culture problem. \\[4pt]
		\textbf{A wrong:} Structure alone is insufficient. \\[2pt]
		\textbf{C wrong:} Zero reports may indicate fear, not ethics. \\[2pt]
		\textbf{D wrong:} Surveys provide relevant culture evidence.
	\end{frame}
	
	% ============================================================
	% Q103
	% ============================================================
	\begin{frame}[t]{Q103}
		\small
		\textbf{Question:} The CAE is asked to evaluate whether the organization's board has effective oversight of risk. The board receives a quarterly risk report that is prepared by management and contains only risk status ("green") with no details on emerging risks or risk appetite breaches. Which conclusion is MOST appropriate?
		\begin{enumerate}[label=\Alph*.,itemsep=1pt,topsep=2pt]
			\item Board risk oversight is effective because the board receives a quarterly report.
			\item Board risk oversight is likely ineffective; the report lacks detail, emerging risks, and appetite breach information.
			\item Risk reporting is management's responsibility, not the board's.
			\item The auditor should not assess board risk oversight.
		\end{enumerate}
	\end{frame}
	
	\begin{frame}[t]{Q103 -- Answer \& Explanation}
		\small
		\textbf{Correct Answer: B} \\[4pt]
		\textbf{Why:} Effective board risk oversight requires transparent, detailed reporting on emerging risks and appetite breaches. A "green-only" summary impairs oversight. \\[4pt]
		\textbf{A wrong:} Receiving a report is not the same as effective oversight. \\[2pt]
		\textbf{C wrong:} The board is responsible for oversight. \\[2pt]
		\textbf{D wrong:} Internal audit assesses governance/risk oversight.
	\end{frame}
	
	% ============================================================
	% Q104
	% ============================================================
	\begin{frame}[t]{Q104}
		\small
		\textbf{Question:} An auditor is testing the organization's controls over the sales commission process. The auditor finds that sales managers can adjust commission rates and approve their own adjustments, with no independent review. Which conclusion is MOST appropriate?
		\begin{enumerate}[label=\Alph*.,itemsep=1pt,topsep=2pt]
			\item Controls are effective because sales managers understand commissions.
			\item Controls are deficient; sales managers can adjust rates and approve their own changes, creating a significant fraud risk.
			\item Controls are effective because commission adjustments are rare.
			\item The auditor should not assess sales commissions.
		\end{enumerate}
	\end{frame}
	
	\begin{frame}[t]{Q104 -- Answer \& Explanation}
		\small
		\textbf{Correct Answer: B} \\[4pt]
		\textbf{Why:} The ability to adjust rates and self-approve creates a significant fraud risk (self-dealing). Independent review is required. \\[4pt]
		\textbf{A wrong:} Knowledge is not a control. \\[2pt]
		\textbf{C wrong:} Rarity does not reduce the fraud risk. \\[2pt]
		\textbf{D wrong:} Commissions are a core audit area.
	\end{frame}
	
	% ============================================================
	% Q105
	% ============================================================
	\begin{frame}[t]{Q105}
		\small
		\textbf{Question:} The CAE is evaluating the organization's compliance with its own code of conduct. The CAE finds that senior executives routinely accept expensive gifts from vendors and that no one has been disciplined. Which conclusion is MOST appropriate?
		\begin{enumerate}[label=\Alph*.,itemsep=1pt,topsep=2pt]
			\item The code of conduct is effective because it is written.
			\item The code of conduct is not effective; senior executive behavior and lack of enforcement undermine the code.
			\item Gifts are acceptable if they are expensive.
			\item The auditor should ignore executive behavior.
		\end{enumerate}
	\end{frame}
	
	\begin{frame}[t]{Q105 -- Answer \& Explanation}
		\small
		\textbf{Correct Answer: B} \\[4pt]
		\textbf{Why:} Tone at the top and enforcement are essential. Senior executives accepting gifts without consequences undermine the code's effectiveness. \\[4pt]
		\textbf{A wrong:} A written code without enforcement is ineffective. \\[2pt]
		\textbf{C wrong:} Expensive gifts create significant conflicts. \\[2pt]
		\textbf{D wrong:} Executive behavior is a key governance indicator.
	\end{frame}
	
	% ============================================================
	% Q106
	% ============================================================
	\begin{frame}[t]{Q106}
		\small
		\textbf{Question:} An auditor is assessing the organization's controls over the use of spreadsheets in financial reporting. The auditor finds that a critical revenue calculation is performed in a spreadsheet with no version control, no formula documentation, and no independent review. Which conclusion is MOST appropriate?
		\begin{enumerate}[label=\Alph*.,itemsep=1pt,topsep=2pt]
			\item Controls are adequate because spreadsheets are common tools.
			\item Controls are deficient; critical spreadsheets require version control, documentation, and independent review.
			\item Spreadsheets are always inappropriate for financial reporting.
			\item The auditor should not assess spreadsheet controls.
		\end{enumerate}
	\end{frame}
	
	\begin{frame}[t]{Q106 -- Answer \& Explanation}
		\small
		\textbf{Correct Answer: B} \\[4pt]
		\textbf{Why:} Critical spreadsheets are high-risk end-user computing tools requiring version control, documentation, and review to prevent errors/fraud. \\[4pt]
		\textbf{A wrong:} Common use does not reduce risk. \\[2pt]
		\textbf{C wrong:} Spreadsheets can be used with adequate controls. \\[2pt]
		\textbf{D wrong:} Spreadsheet controls are a valid audit area.
	\end{frame}
	
	% ============================================================
	% Q107
	% ============================================================
	\begin{frame}[t]{Q107}
		\small
		\textbf{Question:} The CAE is evaluating the organization's risk management for a major IT system implementation. The project has no risk register, no steering committee, and no independent oversight; the project manager reports only to the CIO. Which conclusion is MOST appropriate?
		\begin{enumerate}[label=\Alph*.,itemsep=1pt,topsep=2pt]
			\item Project risk management is adequate because the CIO is accountable.
			\item Project risk management is deficient; major IT projects require a risk register, steering committee, and independent oversight.
			\item IT projects do not require risk management.
			\item The auditor should not assess IT project risk.
		\end{enumerate}
	\end{frame}
	
	\begin{frame}[t]{Q107 -- Answer \& Explanation}
		\small
		\textbf{Correct Answer: B} \\[4pt]
		\textbf{Why:} Major IT projects require structured risk management, governance (steering committee), and independent oversight to manage cost, schedule, and risk. \\[4pt]
		\textbf{A wrong:} CIO accountability alone is insufficient. \\[2pt]
		\textbf{C wrong:} All projects require risk management. \\[2pt]
		\textbf{D wrong:} IT project risk is a core audit area.
	\end{frame}
	
	% ============================================================
	% Q108
	% ============================================================
	\begin{frame}[t]{Q108}
		\small
		\textbf{Question:} An auditor is testing the organization's controls over journal entries. The auditor finds that 20 large manual entries were posted at year-end by the controller with descriptions such as "reclass" and "true-up," and that supporting documentation cannot be located. Which action is MOST appropriate?
		\begin{enumerate}[label=\Alph*.,itemsep=1pt,topsep=2pt]
			\item Accept the entries because the controller is authorized.
			\item Investigate the entries, request support, and report the lack of documentation as a significant deficiency.
			\item Delete the entries from the system.
			\item Ignore the entries because they are year-end adjustments.
		\end{enumerate}
	\end{frame}
	
	\begin{frame}[t]{Q108 -- Answer \& Explanation}
		\small
		\textbf{Correct Answer: B} \\[4pt]
		\textbf{Why:} Large manual year-end entries by senior management are a high fraud risk area. Lack of support is a significant deficiency requiring investigation and reporting. \\[4pt]
		\textbf{A wrong:} Authorization does not eliminate the need for support. \\[2pt]
		\textbf{C wrong:} Internal audit should not alter records. \\[2pt]
		\textbf{D wrong:} Year-end adjustments require documentation.
	\end{frame}
	
	% ============================================================
	% Q109
	% ============================================================
	\begin{frame}[t]{Q109}
		\small
		\textbf{Question:} The CAE is evaluating the organization's use of key risk indicators (KRIs). The organization tracks revenue, profit, and customer satisfaction but has no KRIs for compliance, cybersecurity, or supply chain risk. Which conclusion is MOST appropriate?
		\begin{enumerate}[label=\Alph*.,itemsep=1pt,topsep=2pt]
			\item KRI coverage is adequate because financial and customer metrics are tracked.
			\item KRI coverage is incomplete; the organization lacks KRIs for significant risk areas (compliance, cybersecurity, supply chain).
			\item KRIs are optional for all organizations.
			\item The auditor should not assess KRIs.
		\end{enumerate}
	\end{frame}
	
	\begin{frame}[t]{Q109 -- Answer \& Explanation}
		\small
		\textbf{Correct Answer: B} \\[4pt]
		\textbf{Why:} KRIs should cover the organization's significant risks. Omitting compliance, cybersecurity, and supply chain KRIs leaves key risks unmonitored. \\[4pt]
		\textbf{A wrong:} Financial/customer metrics alone are insufficient. \\[2pt]
		\textbf{C wrong:} KRIs are an important monitoring tool. \\[2pt]
		\textbf{D wrong:} Internal audit assesses risk monitoring.
	\end{frame}
	
	% ============================================================
	% Q110
	% ============================================================
	\begin{frame}[t]{Q110}
		\small
		\textbf{Question:} An auditor is evaluating the organization's controls over the financial close. The auditor finds that the same individual prepares, reviews, and posts consolidation entries, and that the individual has system access to override the consolidation. Which conclusion is MOST appropriate?
		\begin{enumerate}[label=\Alph*.,itemsep=1pt,topsep=2pt]
			\item Controls are effective because the individual is experienced.
			\item Controls are deficient; incompatible duties and override access create a significant risk of undetected error/fraud.
			\item Controls are effective because consolidation is routine.
			\item The auditor should not assess the financial close.
		\end{enumerate}
	\end{frame}
	
	\begin{frame}[t]{Q110 -- Answer \& Explanation}
		\small
		\textbf{Correct Answer: B} \\[4pt]
		\textbf{Why:} Combining prepare/review/post and override access is a severe segregation-of-duties and override risk. \\[4pt]
		\textbf{A wrong:} Experience is not a control. \\[2pt]
		\textbf{C wrong:} Routine processes can still be manipulated. \\[2pt]
		\textbf{D wrong:} Financial close is a core audit area.
	\end{frame}
	
	% ============================================================
	% Q111
	% ============================================================
	\begin{frame}[t]{Q111}
		\small
		\textbf{Question:} The CAE is evaluating the organization's risk appetite for cybersecurity. The board has stated "zero tolerance for data breaches," but management has not funded basic security controls and has accepted known vulnerabilities. Which conclusion is MOST appropriate?
		\begin{enumerate}[label=\Alph*.,itemsep=1pt,topsep=2pt]
			\item Risk appetite and behavior are aligned because the board stated zero tolerance.
			\item There is a disconnect between stated risk appetite and actual behavior; the organization is not resourcing controls consistent with its appetite.
			\item Cybersecurity risk is not a governance issue.
			\item The auditor should ignore the disconnect.
		\end{enumerate}
	\end{frame}
	
	\begin{frame}[t]{Q111 -- Answer \& Explanation}
		\small
		\textbf{Correct Answer: B} \\[4pt]
		\textbf{Why:} Risk appetite must be supported by resources and actions. A stated zero tolerance without funding/controls indicates misalignment and a weak risk culture. \\[4pt]
		\textbf{A wrong:} Statements without action are not alignment. \\[2pt]
		\textbf{C wrong:} Cybersecurity is a governance risk. \\[2pt]
		\textbf{D wrong:} Misalignment is a reportable governance weakness.
	\end{frame}
	
	% ============================================================
	% Q112
	% ============================================================
	\begin{frame}[t]{Q112}
		\small
		\textbf{Question:} An auditor is testing the organization's controls over the use of artificial intelligence (AI) in credit decisions. The auditor finds that the AI model was developed by an external vendor, has not been validated, and that management cannot explain the model's logic. Which conclusion is MOST appropriate?
		\begin{enumerate}[label=\Alph*.,itemsep=1pt,topsep=2pt]
			\item Controls are adequate because the model was developed by a vendor.
			\item Controls are deficient; unvalidated AI models with unexplainable logic create significant model, compliance, and reputational risk.
			\item AI models do not require validation.
			\item The auditor should not assess AI models.
		\end{enumerate}
	\end{frame}
	
	\begin{frame}[t]{Q112 -- Answer \& Explanation}
		\small
		\textbf{Correct Answer: B} \\[4pt]
		\textbf{Why:} AI models used in credit decisions require validation, documentation, and explainability to manage model risk and regulatory compliance. \\[4pt]
		\textbf{A wrong:} Vendor development does not eliminate the need for validation. \\[2pt]
		\textbf{C wrong:} AI models require validation. \\[2pt]
		\textbf{D wrong:} AI risk is a core internal audit area.
	\end{frame}
	
	% ============================================================
	% Q113
	% ============================================================
	\begin{frame}[t]{Q113}
		\small
		\textbf{Question:} The CAE is evaluating the organization's fraud hotline. The hotline is managed by the internal audit function, reports go directly to the CAE, and the CAE investigates all reports. Which statement is MOST accurate?
		\begin{enumerate}[label=\Alph*.,itemsep=1pt,topsep=2pt]
			\item The hotline is effective because internal audit is independent.
			\item The hotline may create an independence/self-review issue if internal audit investigates reports about its own function; oversight by the audit committee is important.
			\item The hotline is ineffective because internal audit manages it.
			\item Hotlines should never be managed by internal audit.
		\end{enumerate}
	\end{frame}
	
	\begin{frame}[t]{Q113 -- Answer \& Explanation}
		\small
		\textbf{Correct Answer: B} \\[4pt]
		\textbf{Why:} Internal audit managing a hotline is common, but reports involving internal audit itself create a self-review threat. Audit committee oversight and escalation protocols are important safeguards. \\[4pt]
		\textbf{A wrong:} Independence alone does not resolve self-review for reports about internal audit. \\[2pt]
		\textbf{C wrong:} Internal audit management is not inherently ineffective. \\[2pt]
		\textbf{D wrong:} Internal audit can manage hotlines with safeguards.
	\end{frame}
	
	% ============================================================
	% Q114
	% ============================================================
	\begin{frame}[t]{Q114}
		\small
		\textbf{Question:} An auditor is testing the organization's controls over the financial close. The auditor finds that the controller prepares, reviews, and posts consolidation entries, and that the controller has system access to override the consolidation. Which conclusion is MOST appropriate?
		\begin{enumerate}[label=\Alph*.,itemsep=1pt,topsep=2pt]
			\item Controls are effective because the controller is experienced.
			\item Controls are deficient; incompatible duties and override access create a significant risk of undetected error/fraud.
			\item Controls are effective because consolidation is routine.
			\item The auditor should not assess the financial close.
		\end{enumerate}
	\end{frame}
	
	\begin{frame}[t]{Q114 -- Answer \& Explanation}
		\small
		\textbf{Correct Answer: B} \\[4pt]
		\textbf{Why:} Combining prepare/review/post and override access is a severe segregation-of-duties and override risk. \\[4pt]
		\textbf{A wrong:} Experience is not a control. \\[2pt]
		\textbf{C wrong:} Routine processes can still be manipulated. \\[2pt]
		\textbf{D wrong:} Financial close is a core audit area.
	\end{frame}
	
	% ============================================================
	% Q115
	% ============================================================
	\begin{frame}[t]{Q115}
		\small
		\textbf{Question:} The CAE is evaluating the organization's risk appetite for cybersecurity. The board has stated "zero tolerance for data breaches," but management has not funded basic security controls and has accepted known vulnerabilities. Which conclusion is MOST appropriate?
		\begin{enumerate}[label=\Alph*.,itemsep=1pt,topsep=2pt]
			\item Risk appetite and behavior are aligned because the board stated zero tolerance.
			\item There is a disconnect between stated risk appetite and actual behavior; the organization is not resourcing controls consistent with its appetite.
			\item Cybersecurity risk is not a governance issue.
			\item The auditor should ignore the disconnect.
		\end{enumerate}
	\end{frame}
	
	\begin{frame}[t]{Q115 -- Answer \& Explanation}
		\small
		\textbf{Correct Answer: B} \\[4pt]
		\textbf{Why:} Risk appetite must be supported by resources and actions. A stated zero tolerance without funding/controls indicates misalignment and a weak risk culture. \\[4pt]
		\textbf{A wrong:} Statements without action are not alignment. \\[2pt]
		\textbf{C wrong:} Cybersecurity is a governance risk. \\[2pt]
		\textbf{D wrong:} Misalignment is a reportable governance weakness.
	\end{frame}
	
	% ============================================================
	% Q116
	% ============================================================
	\begin{frame}[t]{Q116}
		\small
		\textbf{Question:} An auditor is testing the organization's controls over related-party transactions. The auditor finds that the CEO's spouse owns a consulting firm that received \$3 million in contracts, and that the transactions were not disclosed to the board or in the financial statements. Which conclusion is MOST appropriate?
		\begin{enumerate}[label=\Alph*.,itemsep=1pt,topsep=2pt]
			\item The transactions are acceptable because the firm is competitive.
			\item The transactions represent a significant control and disclosure deficiency; related-party transactions must be identified, approved, and disclosed.
			\item Related-party transactions are outside the scope of internal audit.
			\item The auditor should ignore the transactions because they involve a spouse.
		\end{enumerate}
	\end{frame}
	
	\begin{frame}[t]{Q116 -- Answer \& Explanation}
		\small
		\textbf{Correct Answer: B} \\[4pt]
		\textbf{Why:} Related-party transactions require identification, independent approval, and disclosure. Undisclosed transactions with the CEO's spouse's firm are a significant governance/control failure. \\[4pt]
		\textbf{A wrong:} Competitiveness does not eliminate disclosure requirements. \\[2pt]
		\textbf{C wrong:} Related-party transactions are a core audit area. \\[2pt]
		\textbf{D wrong:} Spousal relationships can be fraud indicators.
	\end{frame}
	
	% ============================================================
	% Q117
	% ============================================================
	\begin{frame}[t]{Q117}
		\small
		\textbf{Question:} The CAE is evaluating the organization's cybersecurity governance. The organization has no board-level cybersecurity expertise, no regular cybersecurity reporting to the board, and no incident response plan. Which conclusion is MOST appropriate?
		\begin{enumerate}[label=\Alph*.,itemsep=1pt,topsep=2pt]
			\item Cybersecurity governance is adequate because IT handles security.
			\item Cybersecurity governance is deficient; board oversight, regular reporting, and an incident response plan are essential.
			\item Cybersecurity is purely an IT issue, not a governance issue.
			\item The auditor should not assess cybersecurity.
		\end{enumerate}
	\end{frame}
	
	\begin{frame}[t]{Q117 -- Answer \& Explanation}
		\small
		\textbf{Correct Answer: B} \\[4pt]
		\textbf{Why:} Cybersecurity is a board-level governance risk. Lack of expertise, reporting, and an incident response plan are significant governance deficiencies. \\[4pt]
		\textbf{A wrong:} Security is not solely an IT issue; governance oversight is required. \\[2pt]
		\textbf{C wrong:} Cybersecurity is a governance and enterprise risk. \\[2pt]
		\textbf{D wrong:} Internal audit assesses cybersecurity governance.
	\end{frame}
	
	% ============================================================
	% Q118
	% ============================================================
	\begin{frame}[t]{Q118}
		\small
		\textbf{Question:} An auditor is evaluating the organization's fraud risk assessment. Management's assessment is based solely on a brainstorming session with the CFO and controller, with no input from operations, sales, or IT, and no consideration of fraud schemes specific to the industry. Which conclusion is MOST appropriate?
		\begin{enumerate}[label=\Alph*.,itemsep=1pt,topsep=2pt]
			\item The fraud risk assessment is adequate because senior finance staff participated.
			\item The fraud risk assessment is deficient; it lacks cross-functional input and industry-specific fraud schemes.
			\item The fraud risk assessment is adequate because brainstorming is a recognized technique.
			\item Fraud risk assessment is not an internal audit concern.
		\end{enumerate}
	\end{frame}
	
	\begin{frame}[t]{Q118 -- Answer \& Explanation}
		\small
		\textbf{Correct Answer: B} \\[4pt]
		\textbf{Why:} Effective fraud risk assessment requires cross-functional input and consideration of industry-specific schemes. A finance-only brainstorm is insufficient. \\[4pt]
		\textbf{A wrong:} Finance-only input is too narrow. \\[2pt]
		\textbf{C wrong:} Brainstorming is useful but must be broad and informed. \\[2pt]
		\textbf{D wrong:} Internal audit assesses fraud risk (Standard 1210.A2).
	\end{frame}
	
	\begin{frame}{End}
		\centering
		\Large End of Questions 60--118\\[8pt]
		\normalsize Review, repeat, and mark weak areas.
	\end{frame}
	
\end{document}